\documentclass[11pt]{article}
\usepackage[margin=1in]{geometry}
\usepackage{amsmath,amssymb}
\usepackage{booktabs}
\usepackage{enumitem}
\usepackage{hyperref}
\usepackage{listings}

\title{Chaos Language Algorithm: Technical Summary of\\ Five Standard Attractor System Benchmarks}
\author{ZeroBot / ProtoCosmoBot}
\date{2026-07-07}

\begin{document}
\maketitle

\begin{abstract}
The Chaos Language Algorithm (CLA) converts chaotic dynamical trajectories into symbolic strings and then learns compact grammar productions via n-gram chunk mining, context-category induction, and minimum-description-length (MDL) greedy acceptance. We report benchmark results on five standard attractor systems---Lorenz-63, R\"ossler, Mackey-Glass, Lorenz-96, and the logistic map---with exact reconstruction verified in every case. Dimensionality scaling is explored via Lorenz-96 at dimensions 5, 8, 12, 16, 20, and 24, covering the range from textbook systems to OmegaSim-scale state vectors ($\sim$16 continuous dimensions).
\end{abstract}

\section{Algorithm Overview}

The CLA pipeline consists of four stages:

\begin{enumerate}[noitemsep]
    \item \textbf{Trajectory generation.} Deterministic, fixed-step integrators (RK4 for ODE systems, Euler with discrete delay buffer for Mackey-Glass, direct iteration for the logistic map) produce a sequence of state vectors.
    \item \textbf{M1 symbolization.} Each trajectory point is mapped to a symbol via fixed rectangular partitioning. For an $n$-dimensional trajectory, the symbol is a composite of per-axis bin indices: $\texttt{m1:d0|d1|}\ldots\texttt{|d}_{n-1}$. Bounds are inferred once from the full trajectory and held fixed.
    \item \textbf{Grammar learning.} An n-gram pattern miner discovers repeated substrings of length 2--5. A greedy MDL loop iteratively accepts the chunk or category edit that maximally reduces total description length. Context categories group symbols by co-occurrence using Jensen-Shannon divergence.
    \item \textbf{Verification.} The learned grammar is expanded back to the original symbol stream. If the expansion matches exactly, the run passes.
\end{enumerate}

\section{Benchmark Systems}

\subsection{Lorenz-63}

The classic 3D convection model:
\begin{align}
    \dot{x} &= \sigma(y - x) \\
    \dot{y} &= x(\rho - z) - y \\
    \dot{z} &= xy - \beta z
\end{align}
with canonical parameters $\sigma=10$, $\rho=28$, $\beta=8/3$, initial condition $(1,1,1)$, and $dt=0.01$.

\subsection{R\"ossler}

A 3D system with a simpler attractor structure:
\begin{align}
    \dot{x} &= -y - z \\
    \dot{y} &= x + ay \\
    \dot{z} &= b + z(x - c)
\end{align}
with $a=0.2$, $b=0.2$, $c=5.7$, initial $(0.1, 0, 0)$, and $dt=0.05$.

\subsection{Mackey-Glass}

A scalar delay-differential equation exhibiting chaotic dynamics for large delay:
\begin{equation}
    \dot{x}(t) = \beta \frac{x(t-\tau)}{1 + x(t-\tau)^n} - \gamma x(t)
\end{equation}
with $\beta=0.2$, $\gamma=0.1$, $n=10$, $\tau=17$, initial history $x=0.5$, and $dt=0.1$. Integrated via Euler with a discrete delay buffer of length $\tau$.

\subsection{Lorenz-96}

A variable-dimension system used for dimensionality scaling:
\begin{equation}
    \dot{x}_i = (x_{i+1} - x_{i-2})x_{i-1} - x_i + F, \quad i = 0, \ldots, N-1
\end{equation}
with periodic boundary conditions, $F=8.0$, $dt=0.01$, and initial condition $x_i = 8.0$ for all $i$ except $x_{N-1} = 8.01$. The dimension $N$ is varied from 5 (baseline) through 8, 12, 16, 20, and 24.

\subsection{Logistic Map}

A discrete 1D chaotic map:
\begin{equation}
    x_{n+1} = r x_n(1 - x_n)
\end{equation}
with $r=4.0$ (fully chaotic), $x_0 = 0.123$.

\section{Experimental Configuration}

All experiments use identical CLA parameters:

\begin{center}
\begin{tabular}{ll}
\toprule
Parameter & Value \\
\midrule
Steps & 256 \\
Discard (transient) & 64 \\
Bins per axis & 4 \\
Max iterations & 8 \\
N-gram range & 2--5 \\
Min uses & 2 \\
Categories & Enabled \\
\bottomrule
\end{tabular}
\end{center}

All trajectories are deterministic; no stochastic RNG is used in the integrators. The only randomness is in the CLA grammar-learning seed, which controls tie-breaking in the MDL greedy loop.

\section{Results}

\subsection{Baseline Benchmarks}

\begin{center}
\begin{tabular}{lrrrrr}
\toprule
System & Dim & Unique & Rules & Score (bits) & Bits/sym \\
\midrule
Mackey-Glass & 1 & 4 & 6 & 44.4 & 0.17 \\
Logistic Map & 1 & 4 & 8 & 113.2 & 0.44 \\
Lorenz-63 & 3 & 21 & 8 & 187.2 & 0.73 \\
R\"ossler & 3 & 24 & 8 & 189.2 & 0.74 \\
Lorenz-96 & 5 & 56 & 8 & 234.2 & 0.92 \\
\bottomrule
\end{tabular}
\end{center}

\textbf{Key observations:}
\begin{itemize}[noitemsep]
    \item Exact reconstruction holds for all five systems.
    \item Mackey-Glass achieves the best compression (0.17 bits/symbol) due to its highly repetitive scalar trajectory with only 4 unique symbols.
    \item The logistic map, also 1D, shows moderate compressibility (0.44 bits/symbol) with more diverse symbol sequences.
    \item Lorenz-63 and R\"ossler are closely matched ($\sim$0.74 bits/symbol), consistent with their similar attractor complexity.
    \item Lorenz-96 at dim=5 already shows high symbol diversity (56 unique out of 256), yielding near-logarithmic compression.
\end{itemize}

\subsection{Dimensionality Scaling (Lorenz-96)}

\begin{center}
\begin{tabular}{rrrrrr}
\toprule
Dim & Unique & Possible ($4^N$) & Rules & Score (bits) & Bits/sym \\
\midrule
5 & 56 & 1{,}024 & 8 & 234.2 & 0.92 \\
8 & 78 & 65{,}536 & 8 & 229.2 & 0.90 \\
12 & 107 & $1.7\times10^7$ & 6 & 241.4 & 0.94 \\
16 & 121 & $4.3\times10^9$ & 3 & 248.7 & 0.97 \\
20 & 127 & $1.1\times10^{12}$ & 6 & 251.4 & 0.98 \\
24 & 148 & $2.8\times10^{14}$ & 5 & 251.5 & 0.98 \\
\bottomrule
\end{tabular}
\end{center}

\textbf{Key observations:}
\begin{itemize}[noitemsep]
    \item Exact reconstruction holds at all tested dimensions (1--24).
    \item Unique symbol count grows sub-exponentially relative to the theoretical maximum $4^N$, indicating the trajectory occupies a small fraction of partition space.
    \item Compression degrades monotonically: from 0.92 bits/symbol at dim=5 to 0.98 at dim=24, approaching the entropy of a flat uniform distribution over unique symbols.
    \item At dim=16 (matching OmegaSim's 16-dimensional state vector), CLA discovers only 3 chunk rules and 121 unique symbols, indicating a near-incompressible regime with 4 bins and 256 steps.
    \item At dim=20--24, the system is effectively incompressible at this bin/step resolution. This is the regime where dimension reduction or coarser binning becomes essential.
\end{itemize}

\section{OmegaSim Dimensionality Context}

The OmegaSim A6 thresholded-appraisal model has the following continuous state vector:

\begin{center}
\begin{tabular}{lcc}
\toprule
Component & Dimensions & Range \\
\midrule
Motivation (4 roles) & 4 & $[0, 1]$ \\
Fatigue (4 roles) & 4 & $[0, 1]$ \\
Threshold (4 roles) & 4 & $[0.15, 0.85]$ \\
Artifact maturity & 1 & $[0, 1]$ \\
Provenance debt & 1 & $[0, 1]$ \\
Risk & 1 & $[0, 1]$ \\
Prediction error & 1 & $[0, 1]$ \\
\midrule
\textbf{Total continuous} & \textbf{16} & \\
\bottomrule
\end{tabular}
\end{center}

An additional discrete field (\texttt{last\_action}, $\in \{0,1,2,3\}$) is not symbolized as a continuous variable.

The dimensionality scaling results show that at $N=16$, the M1 symbolization with 4 bins per axis already enters the near-incompressible regime. This has two implications for the CLA$\to$OmegaSim bridge:

\begin{enumerate}
    \item \textbf{Direct M1 symbolization at $N=16$ is feasible but marginal.} CLA maintains exact reconstruction and still discovers a small number of chunk rules, but the grammar captures only the most repetitive subsequences. The symbolic vocabulary (121 unique symbols out of $4^{16} \approx 4.3\times10^9$) is sparse but diverse enough that MDL gains are limited.
    \item \textbf{Dimension reduction or selective symbolization is needed for richer grammar discovery.} Candidate approaches include:
    \begin{itemize}[noitemsep]
        \item Symbolizing only a subset of state variables (e.g., the 4 motivation values, which are the primary behavioral output).
        \item Applying PCA or delay-embedding before symbolization.
        \item Using coarser bins (2--3 per axis) to reduce symbol diversity.
        \item Employing a sliding-window or delay-coordinate symbolization rather than per-timestep partitioning.
    \end{itemize}
\end{enumerate}

\section{Implementation Details}

\subsection{Repository}

\begin{itemize}[noitemsep]
    \item GitHub: \texttt{bgoertzel-sing/chaos-language-algorithm}
    \item Package: \texttt{chaoslang} (pure Python, no external dependencies)
    \item Latest commit: \texttt{847390c} (sprint complete)
    \item Tests: 64 passing
\end{itemize}

\subsection{Sprint Completion Status}

All five sprint tasks are complete:

\begin{center}
\begin{tabular}{cll}
\toprule
Task & Description & Status \\
\midrule
1 & Symbolic-string MVP (chunks, categories, MDL) & \checkmark \\
2 & Recorded benchmark experiments (5 systems) & \checkmark \\
3 & Dimensionality scaling (Lorenz-96 dims 8--24) & \checkmark \\
4 & JS-divergence context clustering & \checkmark \\
5 & Persistence (JSON state, edit-log replay, stable text) & \checkmark \\
\bottomrule
\end{tabular}
\end{center}

\subsection{CLI Usage}

\begin{lstlisting}[frame=single, basicstyle=\small\ttfamily]
# Single system benchmark
PYTHONPATH=src python3 -m chaoslang.benchmarks.m1_controls \
    --system lorenz63 --steps 256 --discard 64 --bins 4 --iterations 8

# Lorenz-96 with custom dimension
PYTHONPATH=src python3 -m chaoslang.benchmarks.m1_controls \
    --system lorenz96 --steps 256 --discard 64 --bins 4 \
    --iterations 8 --dimension 16 --seed 0
\end{lstlisting}

\section{Conclusions and Next Steps}

\begin{enumerate}
    \item CLA successfully discovers grammar structure in all five standard attractor systems with exact reconstruction, validating the core algorithm.
    \item The dimensionality scaling experiments reveal a clear transition from compressible (low-dim) to near-incompressible (dim $\geq 16$) regimes at the current M1 resolution of 4 bins and 256 steps.
    \item The OmegaSim A6 model's 16-dimensional state vector sits exactly at this transition boundary, making it the natural test case for bridging textbook benchmarks to simulation-scale traces.
    \item The next step is to select toy attractor systems at $N \approx 16$ and test dimension-reduction strategies (selective symbolization, PCA, delay-embedding, coarser binning) to improve grammar discovery in this regime.
    \item Once CLA can robustly extract grammar from $N \approx 16$ systems, it can be applied to actual OmegaSim traces to assess whether the thresholded-appraisal dynamics produce complex strange-attractor structure.
\end{enumerate}

\end{document}
